% ============================================================================
% LANGENTWURF: Las asignaturas (Schulfächer)
% ============================================================================
% Sequenz 4: Mi colegio | Block a | Doppelstunde 1-2
% Fachfarbe: #c41e3a (Karminrot)
% Tier 1 (vollständig)
% ============================================================================

\documentclass[11pt,a4paper]{article}

% ============================================================================
% PACKAGES
% ============================================================================
\usepackage[utf8]{inputenc}
\usepackage[T1]{fontenc}
\usepackage[ngerman]{babel}
\usepackage{geometry}
\usepackage{fancyhdr}
\usepackage{lastpage}
\usepackage{xcolor}
\usepackage{tcolorbox}
\usepackage{enumitem}
\usepackage{tabularx}
\usepackage{longtable}
\usepackage{array}
\usepackage{booktabs}
\usepackage{hyperref}
\usepackage{ifthen}
\usepackage{amssymb}

% ============================================================================
% KONFIGURATION
% ============================================================================

\newcommand{\plantier}{1}
\newcommand{\fachid}{spanisch}
\newcommand{\fach}{Spanisch}
\newcommand{\fachfarbe}{c41e3a}

% ============================================================================
% VARIABLEN
% ============================================================================

\newcommand{\jahrgangsstufe}{7}
\newcommand{\schulart}{Gesamtschule}
\newcommand{\stundenthema}{Las asignaturas (Schulfächer)}
\newcommand{\reihenthema}{Mi colegio}
\newcommand{\stundennummer}{4a}
\newcommand{\reihenstunden}{4}
\newcommand{\unterrichtsdatum}{}
\newcommand{\unterrichtszeit}{}
\newcommand{\raum}{}

\newcommand{\lehrkraft}{N.N.}
\newcommand{\ausbildungsstand}{Lehrkraft}

\newcommand{\lerngruppengroesse}{24}
\newcommand{\maennlich}{12}
\newcommand{\weiblich}{12}
\newcommand{\divers}{0}

\newcommand{\gerniveau}{A1}
\newcommand{\lehrwerk}{Que pasa 1}
\newcommand{\lehrwerkseite}{S. 36-42}

% ============================================================================
% LAYOUT
% ============================================================================
\geometry{left=2.5cm, right=2.5cm, top=2.5cm, bottom=2.5cm}
\definecolor{fach}{HTML}{\fachfarbe}
\definecolor{gray}{HTML}{666666}
\definecolor{lightgray}{HTML}{f5f5f5}
\definecolor{successgreen}{HTML}{2a7a4b}
\definecolor{warningorange}{HTML}{b35c00}

\tcbuselibrary{skins,breakable}

\newtcolorbox{infobox}[1][]{
    colback=lightgray,
    colframe=fach,
    fonttitle=\bfseries,
    title=#1,
    breakable
}

\newtcolorbox{scorebox}{
    colback=white,
    colframe=fach,
    fonttitle=\bfseries,
    title=Didaktik-Score,
    breakable
}

\pagestyle{fancy}
\fancyhf{}
\fancyhead[L]{\small\textcolor{gray}{\fach{} -- Jg. \jahrgangsstufe{} -- \stundenthema}}
\fancyhead[R]{\small\textcolor{gray}{Langentwurf Tier 1}}
\fancyfoot[C]{\small\thepage{} / \pageref{LastPage}}
\renewcommand{\headrulewidth}{0.4pt}
\renewcommand{\footrulewidth}{0pt}

% ============================================================================
% DOKUMENT
% ============================================================================
\begin{document}

% ============================================================================
% DECKBLATT
% ============================================================================
\begin{titlepage}
    \centering
    \vspace*{2cm}

    {\large\textcolor{gray}{\schulart}}\\[2cm]

    {\Huge\textcolor{fach}{\textbf{Langentwurf}}}\\[0.5cm]
    {\large Tier 1 -- Vollständig}\\[2cm]

    {\LARGE\textbf{\stundenthema}}\\[0.5cm]
    {\large im Rahmen der Unterrichtsreihe}\\[0.3cm]
    {\Large\textit{\reihenthema}}\\[2cm]

    \begin{tabular}{rl}
        \textbf{Fach:} & \fach{} (\gerniveau) \\[0.3cm]
        \textbf{Klasse:} & \jahrgangsstufe \\[0.3cm]
        \textbf{Lehrwerk:} & \lehrwerk, \lehrwerkseite \\[0.3cm]
        \textbf{Datum:} & \underline{\hspace{3cm}} \\[0.3cm]
        \textbf{Zeit:} & \underline{\hspace{3cm}} (Doppelstunde) \\[0.3cm]
        \textbf{Raum:} & \underline{\hspace{3cm}} \\[0.3cm]
    \end{tabular}

    \vfill

    {\large Lehrkraft: \lehrkraft}\\[0.3cm]
    {\normalsize\textcolor{gray}{\ausbildungsstand}}\\[1cm]

\end{titlepage}

% ============================================================================
% INHALTSVERZEICHNIS
% ============================================================================
\tableofcontents
\newpage

% ============================================================================
% 1. BEDINGUNGSANALYSE
% ============================================================================
\section{Bedingungsanalyse}

\subsection{Lerngruppe}
\begin{tabular}{ll}
    Kursgröße: & \lerngruppengroesse{} SuS (\maennlich{} m / \weiblich{} w / \divers{} d) \\
    Jahrgangsstufe: & \jahrgangsstufe \\
    Sprachniveau: & \gerniveau{} (GER) \\
    Fremdsprache: & 2. Fremdsprache \\
\end{tabular}

\subsection{Differenzierung (intern)}

\textbf{WICHTIG:} Keine Sternchen auf Schülermaterial!

\begin{tabular}{|l|p{4cm}|p{4cm}|p{4cm}|}
\hline
& \textbf{[B] Basis} & \textbf{[S] Standard} & \textbf{[E] Erweiterung} \\
\hline
\textbf{Ziel} & Berufsreife & Mittlere Reife & Gymnasium \\
\hline
\textbf{Inhalt} & Rezeptiv, Zuordnung & Produktion mit Hilfe & Freie Meinungsäußerung \\
\hline
\textbf{Hilfen} & Bildkarten, Wortliste & Beispielsätze & Nur Impulse \\
\hline
\end{tabular}

\subsection{Besonderheiten}
\begin{itemize}
    \item \textbf{LRS:} 2-3 SuS (Nachteilsausgleich: Zeitzugabe, vergrößerte Schrift)
    \item \textbf{DaZ:} 1-2 SuS (Wortschatz deutsch-spanisch vorab besprechen)
    \item \textbf{Leistungsstark:} 3-4 SuS (Begründungen für Meinungen, Zusatzaufgaben)
\end{itemize}

% ============================================================================
% 2. SACHANALYSE
% ============================================================================
\section{Sachanalyse}

\subsection{Sprachliche Strukturen}

\textbf{Wortschatz -- Schulfächer (las asignaturas):}
\begin{center}
\begin{tabular}{|l|l|}
\hline
\textbf{Spanisch} & \textbf{Deutsch} \\
\hline
las matematicas & Mathematik \\
el espanol & Spanisch \\
el ingles & Englisch \\
el aleman & Deutsch \\
la historia & Geschichte \\
la geografia & Erdkunde \\
la biologia & Biologie \\
la musica & Musik \\
el arte & Kunst \\
la educacion fisica & Sport \\
la informatica & Informatik \\
\hline
\end{tabular}
\end{center}

\textbf{Grammatik -- Meinungsäußerung:}
\begin{itemize}
    \item \textbf{Me gusta...} + Singular-Nomen (Me gusta la musica.)
    \item \textbf{Me gustan...} + Plural-Nomen (Me gustan las matematicas.)
    \item \textbf{No me gusta...} / \textbf{No me gustan...} (Verneinung)
    \item \textbf{Porque...} (Begründung: porque es interesante/aburrido/facil/dificil)
\end{itemize}

\textbf{Phonetik:}
\begin{itemize}
    \item Aussprache: ñ in español, ll in llamas
    \item Betonung: matematicas (auf dem zweiten a)
    \item Stimmhafte Konsonanten: biologia, geografia
\end{itemize}

\subsection{Kontrastive Betrachtung (Deutsch $\leftrightarrow$ Spanisch)}
\begin{itemize}
    \item \textbf{Positive Transfers:} Internationale Wörter (matematicas $\approx$ Mathematik, biologia $\approx$ Biologie)
    \item \textbf{Interferenzen:}
    \begin{itemize}
        \item Genus unterschiedlich (la biologia $\neq$ die Biologie)
        \item Gustar-Konstruktion (indirektes Objekt + Verb, nicht Subjekt + Verb)
    \end{itemize}
    \item \textbf{Falsche Freunde:} gimnasio = Fitnessstudio (nicht Gymnasium)
\end{itemize}

\subsection{Kulturelle Aspekte}
\begin{itemize}
    \item Schulsystem in Spanien: ESO (Educacion Secundaria Obligatoria, 12-16 Jahre)
    \item Stundenplan: Unterricht oft von 8:00-14:30, dann Mittagspause
    \item Fächerbezeichnungen können regional variieren (z.B. matematica vs. matematicas)
\end{itemize}

% ============================================================================
% 3. DIDAKTISCHE ANALYSE
% ============================================================================
\section{Didaktische Analyse}

\subsection{Stellung in der Sequenz}
Dies ist die \textbf{\stundennummer. Doppelstunde} der Sequenz ``\reihenthema'' (\reihenstunden{} Doppelstunden gesamt).

\begin{tabular}{|c|l|l|}
\hline
\textbf{Block} & \textbf{Thema} & \textbf{Status} \\
\hline
4a & Las asignaturas (Schulfächer) & \textbf{aktuell} \\
4b & El horario (Stundenplan) & folgt \\
4c & Verbos en -ar (Verben auf -ar) & folgt \\
4d & Mi colegio -- Sequenztest & folgt \\
\hline
\end{tabular}

\subsection{Kompetenzziele (Rahmenplan MV)}
\begin{itemize}
    \item \textbf{Kommunikativ -- Sprechen:} Über Schulfächer sprechen, Vorlieben ausdrücken
    \item \textbf{Kommunikativ -- Hören:} Schulfächer in authentischen Kontexten verstehen
    \item \textbf{Kommunikativ -- Lesen:} Stundenpläne verstehen
    \item \textbf{Sprachlich:} Wortfeld Schulfächer, Struktur gustar + Infinitiv/Nomen
    \item \textbf{Interkulturell:} Spanisches Schulsystem kennenlernen
    \item \textbf{Methodisch:} Wortschatzstrategien (Internationalismen erkennen)
\end{itemize}

\subsection{Legitimation}
\textbf{Rahmenplan Spanisch MV (Kl. 7-10):}
\begin{itemize}
    \item Kompetenzbereich: Kommunikative Kompetenz -- Sprechen
    \item Standard: ``SuS können einfache Aussagen über den eigenen Alltag machen''
    \item Themenfeld: yo y mi entorno -- la escuela
\end{itemize}

% ============================================================================
% 4. STUNDENZIELE
% ============================================================================
\section{Stundenziele}

\subsection{Grobziel}
Die SuS erweitern ihre \textbf{kommunikative Kompetenz} im Bereich Sprechen, indem sie Schulfächer auf Spanisch benennen und ihre Meinung dazu äußern können.

\subsection{Feinziele (operationalisiert)}
\begin{tabular}{|c|p{8cm}|c|c|}
\hline
\textbf{Nr.} & \textbf{Die SuS können...} & \textbf{AFB} & \textbf{Diff.} \\
\hline
FZ1 & mindestens 8 Schulfächer auf Spanisch korrekt aussprechen und dem deutschen Äquivalent zuordnen & I & alle \\
\hline
FZ2 & den bestimmten Artikel (el/la/los/las) zu den Schulfächern korrekt verwenden & I & alle \\
\hline
FZ3 & mit ``Me gusta(n)'' / ``No me gusta(n)'' ihre Meinung zu Schulfächern ausdrücken & II & alle \\
\hline
FZ4 & einen einfachen Stundenplan lesen und Informationen daraus entnehmen & II & [S]/[E] \\
\hline
FZ5 & ihre Meinung mit ``porque + Adjektiv'' begründen & III & [E] \\
\hline
\end{tabular}

% ============================================================================
% 5. METHODISCHE ANALYSE
% ============================================================================
\section{Methodische Analyse}

\subsection{Phasierung (Doppelstunde = 2 x 45 Min.)}
\begin{enumerate}
    \item \textbf{1. Stunde:}
    \begin{itemize}
        \item Einstieg: Bildimpuls (Stundenplan-Symbol, Bücherstapel)
        \item Erarbeitung: Schulfächer mit Bildkarten einführen
        \item Übung: Chorisches Sprechen, Memory-Spiel
        \item Sicherung: Zuordnungsübung (AB Aufgabe 1)
    \end{itemize}
    \item \textbf{2. Stunde:}
    \begin{itemize}
        \item Wiederholung: Schnelles Vokabelquiz
        \item Erarbeitung: Struktur ``Me gusta(n)'' einführen
        \item Anwendung: Meinungsaustausch in Partnerarbeit
        \item Sicherung: Stundenplan lesen (AB Aufgabe 4)
    \end{itemize}
\end{enumerate}

\subsection{Medien}
\begin{itemize}
    \item Tafel/Whiteboard
    \item Lehrwerk: \lehrwerk, \lehrwerkseite
    \item Arbeitsblatt: AB-04a.pdf (20 Punkte)
    \item Lernpfad: LP-04a.html (optional, digital)
    \item Bildkarten: Schulfächer (11 Stück)
    \item Audio: Que pasa 1, Track 18 (Hörverstehen Schulfächer)
\end{itemize}

% ============================================================================
% 6. VERLAUFSPLANUNG
% ============================================================================
\section{Verlaufsplanung}

\textbf{1. Stunde (45 Min.)}

\begin{longtable}{|p{1cm}|p{1.5cm}|p{4cm}|p{3cm}|p{2cm}|p{2cm}|}
\hline
\textbf{Zeit} & \textbf{Phase} & \textbf{Lehrerhandeln} & \textbf{Schülerhandeln} & \textbf{SF} & \textbf{Medien} \\
\hline
\endhead

0--5 & Einstieg & Begrüßung auf Spanisch, zeigt Bild eines Stundenplans: ``Hoy hablamos de las asignaturas'' & Vermuten Thema, nennen bekannte Fächer & Plenum & Tafel, Bild \\
\hline

5--15 & Input & Präsentiert Bildkarten mit Schulfächern, spricht vor, lässt nachsprechen: ``las matematicas, el espanol...'' & Hören zu, sprechen nach (chorisch und einzeln) & Plenum & Bildkarten \\
\hline

15--25 & Übung 1 & Memory-Spiel in Gruppen: Bild + spanisches Wort zuordnen & Spielen Memory, wiederholen Vokabeln & GA (4er) & Memory-Karten \\
\hline

25--35 & Übung 2 & AB ausgeben, Aufgabe 1 (Zuordnung) erklären, unterstützt bei Fragen & Bearbeiten Zuordnung selbstständig & EA & AB-04a \\
\hline

35--40 & Sicherung & Lösungen gemeinsam besprechen, korrigieren, Aussprache üben & Vergleichen, korrigieren, fragen nach & Plenum & Tafel \\
\hline

40--45 & Überleitung & Sammelt Artikel an Tafel (el/la), Ausblick 2. Stunde: ``Meinungen äußern'' & Notieren Artikel im Merkkasten & Plenum & Tafel, AB \\
\hline
\end{longtable}

\vspace{1em}

\textbf{2. Stunde (45 Min.)}

\begin{longtable}{|p{1cm}|p{1.5cm}|p{4cm}|p{3cm}|p{2cm}|p{2cm}|}
\hline
\textbf{Zeit} & \textbf{Phase} & \textbf{Lehrerhandeln} & \textbf{Schülerhandeln} & \textbf{SF} & \textbf{Medien} \\
\hline
\endhead

0--5 & Wiederholung & Schnelles Vokabelquiz: L zeigt Bild, SuS nennen spanisches Wort & Antworten schnell, üben Aussprache & Plenum & Bildkarten \\
\hline

5--15 & Input & Führt ``Me gusta...'' ein mit Beispielen: ``Me gusta la musica. No me gusta la historia.'' Erklärt Singular/Plural & Hören zu, notieren Struktur, fragen nach & Plenum & Tafel \\
\hline

15--25 & Übung 3 & Aufgabe 3 (Meinung äußern) erklären, Satzbausteine an Tafel & Schreiben eigene Sätze mit Me gusta(n) & EA & AB-04a \\
\hline

25--35 & Anwendung & Partnerarbeit: ``Frage deinen Partner nach 3 Fächern'' Modelliert mit Schüler/in & Führen Mini-Dialoge: ``?Te gusta el espanol?'' -- ``Si, me gusta.'' & PA & -- \\
\hline

35--42 & Festigung & Aufgabe 4 (Stundenplan lesen) bearbeiten lassen & Entnehmen Informationen aus Stundenplan & EA & AB-04a \\
\hline

42--45 & Abschluss & Zusammenfassung, HA ansagen: Vokabeln lernen, Lernpfad optional & Notieren HA, räumen auf & Plenum & -- \\
\hline
\end{longtable}

\textbf{Legende:} EA = Einzelarbeit, PA = Partnerarbeit, GA = Gruppenarbeit, SF = Sozialform

% ============================================================================
% 7. ERWARTETE SCHWIERIGKEITEN
% ============================================================================
\section{Erwartete Schwierigkeiten}

\begin{tabular}{|p{5cm}|p{8cm}|}
\hline
\textbf{Schwierigkeit} & \textbf{Reaktion} \\
\hline
Aussprache: matematicas (Betonung) & Mehrfach vorsprechen, Silben klatschen: ma-te-MA-ti-cas \\
\hline
Verwechslung: gusta vs. gustan & Visualisierung: Ein Fach = gusta, mehrere = gustan; Farben nutzen \\
\hline
Artikel vergessen & Immer mit Artikel lernen lassen, Merkregel: ``Fach + Artikel = Paar'' \\
\hline
Interferenz: ``Ich mag'' statt Gustar-Struktur & Bewusstmachung: ``Das Fach gefällt mir'' (nicht ``Ich mag das Fach'') \\
\hline
Zeitproblem & Puffer: Aufgabe 4 als HA auslagern \\
\hline
Heterogenität & [B]: Mit Bildkarten unterstützen, [E]: Begründungen einfordern \\
\hline
\end{tabular}

% ============================================================================
% 8. TAFELBILD
% ============================================================================
\section{Tafelbild}

\begin{center}
\fbox{
\begin{minipage}{0.9\textwidth}
\centering
\textbf{\Large Las asignaturas} \\[1em]

\begin{tabular}{|c|c||c|c|}
\hline
\textbf{el} (maskulin) & & \textbf{la} (feminin) & \\
\hline
el espanol & Spanisch & la historia & Geschichte \\
el ingles & Englisch & la geografia & Erdkunde \\
el aleman & Deutsch & la biologia & Biologie \\
el arte & Kunst & la musica & Musik \\
& & la educacion fisica & Sport \\
& & la informatica & Informatik \\
\hline
\multicolumn{4}{|c|}{\textbf{las matematicas} (immer Plural!)} \\
\hline
\end{tabular}

\vspace{1em}

\textbf{Meinung ausdrücken:} \\
\textcolor{green!60!black}{Me gusta} + \textit{el/la + Fach (Singular)} \\
\textcolor{green!60!black}{Me gustan} + \textit{las matematicas (Plural)} \\
\textcolor{red!70!black}{No me gusta(n)} + ... \\[0.5em]
\textit{Porque es...} interesante / aburrido / facil / dificil

\end{minipage}
}
\end{center}

% ============================================================================
% 9. HAUSAUFGABE
% ============================================================================
\section{Hausaufgabe}

\begin{itemize}
    \item Vokabeln lernen: Schulfächer (11 Wörter mit Artikel)
    \item Lehrwerk S. 38, Aufgabe 3 (schriftlich)
    \item Optional: Lernpfad LP-04a.html digital bearbeiten
\end{itemize}

% ============================================================================
% 10. ANLAGEN
% ============================================================================
\section{Anlagen}

\begin{enumerate}
    \item Arbeitsblatt (AB-04a.pdf) -- 20 Punkte
    \item Musterlösung (ML-04a.pdf)
    \item Lehrerhinweise (LH-04a.pdf)
    \item Lernpfad (LP-04a.html)
    \item Bildkarten Schulfächer (11 Stück)
    \item Audio-Track 18 (Que pasa 1)
\end{enumerate}

% ============================================================================
% 11. DIDAKTIK-SCORE
% ============================================================================
\newpage
\section{Didaktik-Score}

\begin{scorebox}
\textbf{Bewertung der didaktischen Qualität nach 10 Dimensionen}\\
\textbf{Ziel-Score: $\geq$ 9.0 (Premium-Qualität)}

\vspace{1em}
\begin{tabular}{|c|l|c|c|p{4.5cm}|}
\hline
\textbf{\#} & \textbf{Dimension} & \textbf{Gew.} & \textbf{Score} & \textbf{Notiz} \\
\hline
1 & Differenzierung & 15\% & 9/10 & [B] Bildkarten, [S] Satzbausteine, [E] Begründungen \\
\hline
2 & Taxonomie (AFB) & 10\% & 9/10 & AFB I: 40\%, AFB II: 45\%, AFB III: 15\% \\
\hline
3 & Lernziel-Alignment & 10\% & 10/10 & Exakt auf RP MV abgestimmt \\
\hline
4 & Aktivierung & 10\% & 9/10 & Memory, PA-Dialoge, eigenständige Satzproduktion \\
\hline
5 & Scaffolding & 10\% & 9/10 & Gestufte Hilfen, Beispielsätze, Strukturhilfen \\
\hline
6 & Feedback-Qualität & 10\% & 9/10 & Sofortige Korrektur, elaborierte Rückmeldung \\
\hline
7 & Sprachliche Klarheit & 10\% & 10/10 & Altersgerecht, klare Operatoren \\
\hline
8 & Motivation \& Relevanz & 10\% & 9/10 & Alltagsbezug (eigener Stundenplan), Spielelement \\
\hline
9 & Inklusion (UDL) & 10\% & 9/10 & Visuell (Bilder), auditiv (Hören), kinästhetisch (Memory) \\
\hline
10 & Praktikabilität & 5\% & 9/10 & Realistisch in 90 Min., Material vorhanden \\
\hline
\multicolumn{3}{|r|}{\textbf{GESAMT:}} & \textbf{9.15/10} & \\
\hline
\end{tabular}

\vspace{1em}
\textbf{Status:} \textcolor{successgreen}{$\checkmark$ FREIGABE (Premium-Qualität)}

\vspace{1em}
\textbf{Stärken:}
\begin{itemize}
    \item Klare Progression: Rezeption $\rightarrow$ Reproduktion $\rightarrow$ Produktion
    \item Abwechslungsreiche Methoden (Memory, PA-Dialoge, Stundenplan)
    \item Internationalismen als Lernstrategie genutzt
    \item Gustar-Struktur didaktisch reduziert auf wesentliche Formen
\end{itemize}

\textbf{Hinweise zur Optimierung:}
\begin{itemize}
    \item Bei Zeitmangel: Aufgabe 4 als HA
    \item Für schwächere SuS: Mehr Zeit für chorisches Sprechen
    \item Für stärkere SuS: Begründungen mit weil einfordern
\end{itemize}

\end{scorebox}

\end{document}
